\chapter{Persuasión, Influencia y Negociación}\label{ch:persuasion-influence}

% Alcance: influencia basada en evidencia --- principios validados por experimentos de campo,
%   encuadre de procesamiento dual, negociación basada en principios (BATNA/valor de reserva/ZOPA),
%   poder organizacional, la ecuación de confianza. NUEVO (Parte VII).
% Gancho de valor: la influencia convierte una propuesta de valor genuina en clientes cerrados
%   (CAC menor), rondas financiadas, contrataciones y excedente capturado (margen mayor) ---
%   todo orientado hacia ROIC > WACC. La manipulación lo invierte (ganancia a corto plazo,
%   colapso de confianza/LTV). Cierra la brecha de negociación.
% Borrador conforme al estilo editorial: prosa continua + example/admonition; ejemplos trabajados
%   con la SaaS de referencia.

Un producto que genera valor genuino aún puede fracasar en capturarlo sin la capacidad de
persuadir. La columna vertebral de creación de valor en este libro va desde la comprensión
económica hasta la ejecución operativa, pero cada nodo de esa columna --- cerrar un cliente,
financiar una ronda, contratar a un ingeniero, fijar el precio con un proveedor --- requiere
que una persona decida moverse en la dirección correcta. La influencia es el mecanismo que
convierte una propuesta de valor genuina en una venta cerrada (reduciendo el CAC, como en
\cref{eq:cac}), una hoja de términos firmada y excedente capturado (ampliando el margen y
empujando el ROIC por encima del WACC, según \cref{eq:roic}). Cuando las técnicas de
influencia se usan en cambio para fabricar una venta que el producto subyacente no puede
sostener, la conversión a corto plazo se revierte en el mediano plazo a través del \emph{churn}
que colapsa el valor de vida del cliente (véase \cref{eq:clv}). Este capítulo construye el
instrumental basado en evidencia que separa la persuasión legítima de la manipulación, y la
negociación basada en principios del regateo posicional de suma cero.


\section{La arquitectura de la influencia}\label{sec:influence-principles}
% complejidad: 3

¿Por qué algunos mensajes mueven a las personas mientras otros se disuelven en el ruido? La
pregunta tiene una respuesta estructural fundamentada en experimentos de campo, no en
anécdotas. El tratamiento popular de este dominio, \textcite{cialdini2021influence}, cataloga
siete principios de cumplimiento; el fundamento científico detrás de esos principios es un
conjunto de experimentos de campo aleatorizados que miden el comportamiento de cumplimiento
real en entornos reales, con condiciones de control preasignadas y tamaños del efecto
cuantificados.

La evidencia experimental más directa concierne a las normas descriptivas --- la señal de
«las personas como tú, en tu situación, se comportan de esta manera». \textcite{goldstein2008field}
llevaron a cabo un experimento de campo dentro de un hotel, aleatorizando los mensajes de
reutilización de toallas entre habitaciones ocupadas sin que los huéspedes lo supieran. El
mensaje ambiental estándar («Ayude a salvar el medio ambiente») produjo una tasa de
reutilización de \qty{35.1}{\percent}. Una norma descriptiva general («La mayoría de los
huéspedes reutilizan sus toallas») elevó el cumplimiento hasta \qty{44.1}{\percent} --- un
incremento absoluto de \qty{9}{\percent} puntos, que representa un aumento relativo de
aproximadamente \qty{25.6}{\percent} sobre el control.
% goldstein2008field: aumento por norma descriptiva vs. control, experimento de campo con toallas de hotel
Una segunda condición ajustó la norma al contexto inmediato del individuo: «El \qty{75}{\percent}
de los huéspedes que se alojaron en esta habitación reutilizaron sus toallas.» Esta norma
provincial elevó el cumplimiento hasta \qty{49.3}{\percent}.
% goldstein2008field: tasa de norma provincial, segundo experimento
El mecanismo es el mismo que rige la optimización de la tasa de conversión: las personas usan
el comportamiento observado de un grupo de referencia para resolver la incertidumbre sobre la
acción correcta, y cuanto más de cerca ese grupo de referencia refleja su propia situación,
mayor es la atracción.

La reciprocidad ha sido medida con igual rigor en el campo. \textcite{strohmetz2002field}
aleatorizaron la manera en que los meseros de un restaurante entregaban la cuenta final,
introduciendo pequeños obsequios inesperados para medir el incremento en propinas. Un solo
caramelo elevó las propinas un \qty{3.3}{\percent} con respecto al control sin obsequio; dos
caramelos las elevaron un \qty{14.1}{\percent}.
% strohmetz2002field: aumento con un solo caramelo 3.3%, con dos caramelos 14.1%, experimento de campo en propinas
El efecto más pronunciado --- un aumento relativo de aproximadamente \qty{23}{\percent} ---
ocurrió cuando el mesero entregó un caramelo, se alejó y luego regresó para ofrecer un segundo
como excepción personalizada.
% strohmetz2002field: secuencia personalizada, pico ~23% de aumento relativo
Dos características de este resultado importan para los profesionales. Primera, la magnitud:
el obsequio era trivial en términos monetarios; el incremento en el cumplimiento no lo era.
Segunda, la personalización: la percepción de que la concesión estaba dirigida específicamente
a esa contraparte amplificó el efecto de manera sustancial. Para un movimiento de ventas SaaS,
una extensión de prueba de funcionalidades enmarcada como una excepción exclusiva e inesperada
para el contexto de evaluación específico de ese prospecto activa ambas propiedades de manera
simultánea.

Estos indicadores heurísticos se sitúan en la ruta periférica del Modelo de Probabilidad de
Elaboración (ELM), desarrollado completamente por \textcite{petty1986elaboration}. El ELM
postula que la persuasión recorre una de dos rutas: cuando un comprador tiene alta motivación
y capacidad para pensar con cuidado, sigue la \emph{ruta central} --- sopesando la calidad de
los argumentos, escrutando la evidencia, actualizando las creencias de manera racional. Cuando
la motivación o la capacidad es baja --- distraído, con poco tiempo, poco familiarizado con el
dominio --- sigue la \emph{ruta periférica}, apoyándose en indicadores heurísticos como normas
descriptivas, señales de autoridad o la afinidad que siente por el emisor. El cambio de actitud
por la ruta central es duradero y resistente a la contraargumentación; el cambio de actitud por
la ruta periférica se desvanece cuando el indicador desaparece. \textcite{carpenter2015meta}
sintetizaron 134 estimaciones de tamaño del efecto de la literatura del ELM y confirmaron esta
asimetría de manera cuantitativa: la calidad de los argumentos produce un efecto de persuasión
sustancialmente mayor cuando la elaboración es alta (procesamiento central) que cuando es baja
(procesamiento periférico).
% carpenter2015meta: 134 efectos, efecto de calidad de argumentos en ruta central >> ruta periférica

La implicación práctica de diseño para un \emph{pitch} B2B es la satisfacción simultánea de
ambas rutas: evidencia económica rigurosa para el comprador que profundiza en los indicadores
unitarios, y señales creíbles de prueba social para los primeros treinta segundos antes de que
decida profundizar. Los casos genéricos («miles de equipos nos usan») explotan deficientemente
la prueba social porque la norma está demasiado distante de la situación del prospecto. El
hallazgo de \textcite{goldstein2008field} apunta a una prescripción concreta: segmentar la
prueba social para que coincida con el contexto exacto del prospecto --- industria, etapa de
la empresa, entorno regulatorio --- y el incremento en el cumplimiento se deriva del mecanismo
de norma provincial más que de un halo de autoridad más amplio pero más débil.

Vinculando la influencia con el \emph{funnel} de adquisición (\cref{eq:funnel-chain}): cada
etapa de la cadena --- sesión a prueba, prueba a pago --- es una decisión de cumplimiento.
La prueba social y la reciprocidad son multiplicadores de la tasa de conversión que operan
sobre el denominador del CAC. Una tasa de conversión de prueba a pago del \qty{25}{\percent}
implica que tres de cada cuatro pruebas se pierden sin acción adicional; añadir una señal de
prueba social contextualizada al momento del final de la prueba y un indicador de reciprocidad
genuina puede desplazar esa tasa de manera significativa. Los mismos mecanismos se aplican a
cada puerta del \emph{funnel}.

\begin{example}[Prueba social contextualizada en la puerta de conversión de prueba a pago]
El \emph{funnel} SaaS convierte \num{10000} sesiones en \num{800} pruebas y \num{200} cuentas
de pago por mes, con una tasa de conversión de prueba a pago del \qty{25}{\percent}. El CAC
se sitúa en \money{600} frente a un objetivo de LTV:CAC de \num{5.25}\texttimes{}
(\cref{eq:ltv-cac}).

El equipo añade dos mensajes en el punto de contacto del día 7 de la prueba: (1)~una norma
provincial de un cliente en el segmento exacto del prospecto («El \qty{73}{\percent} de los
equipos SaaS fintech en Series~B de su región convirtieron a pago este trimestre»), activando
el mecanismo de norma provincial de \textcite{goldstein2008field}; (2)~un indicador de
reciprocidad personalizado de un ejecutivo de cuenta con nombre que ofrece una sesión de
\emph{onboarding} extendida, enmarcada como excepción específica para la complejidad de
integración de este prospecto, activando el mecanismo de concesión personalizada de
\textcite{strohmetz2002field}. Suponga una mejora conservadora de \qty{4}{\percent} puntos en
la conversión de prueba a pago, del \qty{25}{\percent} al \qty{29}{\percent}, sustancialmente
por debajo de los topes del experimento de campo. Eso convierte \num{800} pruebas en \num{232}
cuentas de pago --- \num{32} clientes adicionales --- sin cambiar sesiones, pruebas ni el gasto
de adquisición de \money{120000}. El CAC efectivo cae de \money{600} a aproximadamente
\money{517}. Con un CLV de \money{3150} (\cref{eq:clv}), esos \num{32} clientes adicionales
representan \money{100800} en valor presente de capital de cliente creado sin desembolso
marginal alguno. La palanca es el encuadre en una puerta de cuello de botella, no gasto
adicional en medios.
\end{example}

\begin{admonition}{Advertencia}
Los mismos mecanismos que permiten la persuasión ética pueden ser utilizados como armas. La
escasez falsa («solo quedan 3» cuando el estante está lleno), la prueba social fabricada
(reseñas compradas, conteos de usuarios inflados) o la urgencia manufacturada son patrones
oscuros: elevan la métrica de conversión inmediata y colapsan las dos cosas que la sostienen.
La confianza se evapora ante la primera contradicción, generando \emph{churn} que destruye
el CLV (\cref{eq:clv}) que la conversión debía crear. El género de literatura de negocios que
trata la manipulación como estrategia --- \textcite{greene1998laws} siendo un ejemplo
prominente --- confunde la victoria táctica a corto plazo con la ventaja competitiva duradera.
La evidencia a largo plazo, analizada en \cref{sec:power-trust}, apunta en la dirección
opuesta: la confianza se acumula en LTV; la manipulación se acumula en \emph{churn} y daño
reputacional.
\end{admonition}


\section{La psicología del pitch}\label{sec:pitch-psychology}
% complejidad: 3

Todo \emph{pitch} es un ejercicio de encuadre: los mismos hechos, presentados de manera
diferente, producen decisiones diferentes. Esto no es un truco retórico --- es la estructura
documentada de la percepción humana del valor. El mecanismo estructural más profundo es la
función de valor asimétrica formalizada en la \emph{teoría prospectiva acumulativa} por
\textcite{tversky1992cumulative}. El tratamiento popular de esta comprensión aparece en
\textcite{kahneman2011}; la evidencia cuantitativa primaria es el artículo de 1992, que deriva
el coeficiente de aversión a la pérdida a partir de elecciones experimentales sobre resultados
monetarios reales.

Las personas no evalúan los resultados en magnitudes absolutas; los evalúan como ganancias o
pérdidas relativas a un punto de referencia, y las pérdidas pesan aproximadamente \num{2.25}
veces más que las ganancias equivalentes ($\lambda \approx 2.25$, \cref{eq:prospect}).
% tversky1992cumulative: coeficiente de aversión a la pérdida lambda ~= 2.25, teoría prospectiva acumulativa
Los efectos de contabilidad mental formalizados por \textcite{thaler1985mental} --- y analizados
en el capítulo de precios en el tratamiento del encuadre psicológico en \cref{sec:psych-pricing}
--- son la aplicación directa a las decisiones de compra y pago: la misma cantidad en dólares
que evoca una pérdida se siente peor que la misma cantidad que evoca una ganancia, por lo que
posicionar una oferta como una pérdida evitada en lugar de una ganancia capturada incrementa
la urgencia percibida de actuar. Con $\lambda \approx 2.25$, un \emph{pitch} que enmarca un
beneficio como «está perdiendo actualmente \money{400}/mes» ejerce más del doble de la gravedad
psicológica que un \emph{pitch} que enmarca el beneficio idéntico como «podría ganar
\money{400}/mes».
% tversky1992cumulative: lambda ~= 2.25 implica que el encuadre de pérdida ejerce ~2x la gravedad persuasiva del encuadre de ganancia

La heurística de anclaje, identificada por \textcite{tversky1974judgment}, es el mecanismo
complementario: el primer número en una conversación ejerce una atracción gravitacional sobre
todas las estimaciones posteriores. \textcite{galinsky2001anchoring} demostraron en experimentos
controlados que la primera oferta en una negociación correlaciona con el precio de acuerdo final
en $r = .93$ en condiciones sin defensa, explicando más de la mitad --- y en algunas
especificaciones hasta el \qty{80}{\percent} --- de la varianza en el precio final.
% galinsky2001anchoring: primera oferta / acuerdo final r = .93, hasta 80% de varianza explicada
\textcite{northcraft1987anchoring} extendieron el resultado a expertos del dominio: los agentes
inmobiliarios profesionales a quienes se les dieron diferentes anclas de precio de lista
produjeron tasaciones de propiedades que seguían el ancla, incluso cuando creían que sus
valoraciones se basaban en análisis de mercado objetivo.
% northcraft1987anchoring: anclaje de expertos, profesionales inmobiliarios susceptibles al ancla de precio de lista
En un contexto de \emph{pitch}, el anclaje funciona antes de que comience cualquier negociación:
nombrar el costo de la inacción, o el valor total del resultado, antes de nombrar el precio
establece la referencia contra la cual se juzgará el precio. En un contexto de negociación, la
misma lógica se aplica dentro del ZOPA, tratado en detalle en \cref{sec:negotiation}.

La teoría del procesamiento dual distingue el Sistema~1 (rápido, automático, orientado por
patrones, con peso emocional) del Sistema~2 (lento, deliberado, laborioso, lógico). El
Sistema~1 se activa primero y da forma al marco dentro del cual el Sistema~2 opera
posteriormente. Una presentación de diapositivas que abre con el dolor visceral de un cliente
activa la empatía del Sistema~1 del inversor antes de que su Sistema~2 comience a evaluar los
indicadores unitarios. Invertir el orden --- comenzar con la hoja de cálculo --- obliga al
Sistema~2 a asumir una postura de crítica antes de que exista algún punto de apoyo emocional.
La prescripción práctica es secuenciar: enganchar al Sistema~1 con un problema vívido, luego
satisfacer al Sistema~2 con evidencia rigurosa. El ELM (\cref{sec:influence-principles})
explica por qué ambas rutas deben satisfacerse para que el cambio de actitud persista.

Para un \emph{pitch}, la aversión a la pérdida se traduce en un reencuadre disciplinado de la
propuesta de valor. La pregunta no es «¿qué ganará?» sino «¿qué está perdiendo actualmente,
y por cuánto tiempo?» El punto de referencia es el statu quo, y el statu quo debe costarle
algo visible a la contraparte.

\begin{example}[Encuadrar el nivel empresarial como una pérdida evitada]
La SaaS ofrece un nivel empresarial a \money{299}/mes (frente al nivel estándar a \money{50}/mes).
Bajo un encuadre de ganancia, el \emph{pitch} dice: «Actualice para obtener analítica avanzada,
soporte prioritario y SSO --- un valor de aproximadamente \money{200}/mes para equipos de su
tamaño.» Bajo un encuadre de pérdida anclado al costo del statu quo: «Los equipos de su escala
sin una capa de analítica unificada suelen perder el \qty{15}{\percent} de su eficiencia de
implementación por la conciliación manual --- con su número de empleados, eso equivale a
aproximadamente \money{400} de trabajo improductivo por mes. El nivel empresarial cierra esa
brecha por \money{299}.»

La sustancia es idéntica; el punto de referencia ha cambiado de cero-ganancia a una pérdida
continua. Con $\lambda \approx 2.25$ \parencite{tversky1992cumulative}, la teoría prospectiva
(\cref{eq:prospect}) predice que un comprador en el dominio de pérdidas se vuelve propenso al
riesgo --- más dispuesto a pagar por la certeza de evitar la pérdida que en el dominio de
ganancias por un beneficio equivalente. El nivel de \money{299} ahora tiene un precio inferior
al dolor mensual de \money{400}, enmarcándolo no como un costo sino como una recuperación de
\money{101} por mes. El Sistema~1 cierra el argumento antes de que el Sistema~2 lo vuelva a
abrir. Para respaldar ese cierre del Sistema~2, la ruta central (\cref{sec:influence-principles})
exige evidencia rigurosa de retorno sobre la inversión: un testimonio de un cliente con nombre
con resultados de eficiencia atribuidos y una cifra de tiempo de recuperación que un aprobador
financiero pueda verificar.
\end{example}


\section{Reversión de riesgo y garantías}\label{sec:risk-reversal}
% complejidad: 2

Toda compra es un lanzamiento de moneda que el comprador evalúa: un desembolso cierto ahora
frente a un resultado incierto después. La teoría prospectiva \parencite{tversky1992cumulative}
pone un número concreto en la asimetría --- el coeficiente de aversión a la pérdida
$\lambda \approx 2.25$ significa que una pérdida incierta pesa aproximadamente $2.25 \times$
más en la mente del comprador que una ganancia cierta de igual tamaño. Una garantía traslada
esa pérdida a la columna del vendedor: la desutilidad esperada del comprador por «el producto
no funciona» cae de $\lambda \cdot \text{loss}$ hacia cero, y el vendedor absorbe el riesgo
residual a cambio de una tasa de conversión más alta.

La condición de primer orden es directa. Sea $p$ la fracción en la que la garantía eleva la
tasa de conversión desde su línea base sin garantía, y $\theta$ la tasa de reembolso
postcompra. La garantía es neta positiva siempre que $(1+p)(1-\theta) > 1$, lo que se resuelve
como:

\begin{equation}\label{eq:guarantee-breakeven}
  \theta^{*} \;=\; \frac{p}{1+p}.
\end{equation}

Una garantía que eleva la conversión en \qty{28}{\percent} ($p = 0.28$) soporta tasas de
reembolso de hasta $\theta^{*} = 0.28/1.28 \approx \qty{22}{\percent}$ antes de que la
matemática se vuelva neta negativa --- un margen significativo para absorber el impuesto
de la honestidad.

\begin{example}[Garantía de devolución de dinero en el nivel empresarial]
La SaaS B2B de referencia vende el nivel empresarial a \money{299}/mes, con CLV \money{3150}
(\cref{eq:clv}). La conversión de prueba a pago se sitúa en \qty{25}{\percent} sin garantía;
añadir una garantía de devolución de dinero de \num{90} días sin bloqueo contractual la eleva
al \qty{32}{\percent} ($p = \qty{28}{\percent}$ de aumento en la conversión, en línea con el
multiplicador $\lambda \approx 2.25$ sobre el sacrificio percibido). De los que compran, el
\qty{8}{\percent} solicita el reembolso. Por cada cohorte de 800 pruebas:

\begin{itemize}
  \item Sin garantía: $800 \times 0.25 = 200$ clientes de pago; CLV-VP $= 200 \times \money{3150} = \money{630000}$.
  \item Con garantía: $800 \times 0.32 = 256$ de pago; \num{20} solicitan el reembolso; neto retenido \num{236}; CLV-VP $= 236 \times \money{3150} = \money{743400}$.
\end{itemize}

La garantía añade \money{113400} de VP por cohorte de 800 pruebas a pesar de pagar una tasa
de reembolso del \qty{8}{\percent} --- muy por debajo del umbral de rentabilidad del
\qty{22}{\percent} de \cref{eq:guarantee-breakeven}. La aritmética de la aversión a la pérdida
favorece al vendedor.
\end{example}

\begin{admonition}{Advertencia}
Una garantía invita a la selección adversa: los compradores que sospechan que el producto no
es para ellos son más propensos a solicitar el reembolso, elevando $\theta$ por encima de la
tasa natural. Dos diseños reducen esto: (i)~exigir un hito de uso medible antes de que se
abra el periodo de reembolso (fuerza la autocalificación), y (ii)~acotar el reembolso a la
porción no utilizada (limita el abuso sin eliminar el efecto de reversión de pérdida). La
garantía vacía --- «le encantará» --- no añade ni conversión ni reducción de fricción; la
garantía vinculante, con un mecanismo claro, sí lo hace.
\end{admonition}

La misma lógica se generaliza. Las pruebas gratuitas, los pilotos de pago, los precios de
expansión progresiva y los SLA de riesgo compartido trasladan la incertidumbre de los libros
del comprador a los del vendedor. Cada uno es un candidato cuando el incremento de conversión
que genera supera la pérdida residual que absorbe --- la comparación que formaliza
\cref{eq:guarantee-breakeven}.

\section{Negociación: BATNA, valor de reserva y ZOPA}\label{sec:negotiation}
% complejidad: 4

La persuasión cambia creencias; la negociación asigna el excedente que le sigue. El marco de
negociación basada en principios desarrollado en el Proyecto de Negociación de Harvard ---
cuyo tratamiento popular es \textcite{fisher2011getting} --- descansa en un andamiaje
cuantitativo cuya derivación formal se debe a \textcite{raiffa1982art}. Sus prescripciones
estructurales --- separar a las personas del problema, centrarse en los intereses y no en las
posiciones, inventar opciones para la ganancia mutua, insistir en criterios objetivos ---
son mecanismos para llegar a acuerdos que sean eficientes (sin valor dejado sobre la mesa) y
duraderos (ambas partes perciben el proceso como justo, sosteniendo la relación a través de
transacciones futuras).

Antes de entrar en cualquier negociación, una parte debe responder tres preguntas en privado.
¿Cuál es mi BATNA --- la Mejor Alternativa a un Acuerdo Negociado, el resultado que puedo
obtener sin este acuerdo? ¿Qué implica mi BATNA en términos monetarios como precio de salida?
¿Y dónde se sitúa el límite equivalente de la contraparte? \textcite{kim2005power} establecen
desde la literatura de comportamiento organizacional que el poder estructural en una negociación
se deriva de la fortaleza del BATNA: una parte que puede retirarse de manera creíble mantiene
su posición sin farolear, porque la amenaza es real.
% kim2005power: fortaleza del BATNA como fuente de poder de negociación, explicación estructural
La expresión monetaria del BATNA es el \emph{valor de reserva} --- el punto en o más allá del
cual el acuerdo es peor que ningún acuerdo.

\begin{equation}\label{eq:reservation}
  RV \;=\; U(\mathrm{BATNA}),
\end{equation}
donde \(U(\mathrm{BATNA})\) es la utilidad monetaria de la mejor alternativa externa, completamente
costeada (incluidos los costos de transición, el tiempo y la opcionalidad sacrificada). El valor
de reserva no es una posición de negociación; es un cálculo privado, fijado antes de que comience
la conversación. Revelarlo es casi siempre un error.

La \emph{zona de posible acuerdo} (ZOPA) es el intervalo entre los valores de reserva de
ambas partes en el que ambas prefieren un acuerdo a su opción externa:
\begin{equation}\label{eq:zopa}
  \mathrm{ZOPA} \;=\; [\,RV_{\text{seller}},\; RV_{\text{buyer}}\,],
\end{equation}
con amplitud \(\Delta = RV_{\text{buyer}} - RV_{\text{seller}}\). Si \(\Delta > 0\) un acuerdo
es racionalmente posible y la negociación trata sobre cómo dividir el excedente \(\Delta\). Si
\(\Delta \leq 0\) no existe ningún acuerdo de un solo tema al valor nominal, y las partes deben
introducir nuevas variables (plazos de pago, compromisos de volumen, condiciones de servicio)
para desplazar uno o ambos valores de reserva, o retirarse. Cada táctica --- anclaje, ritmo de
concesiones, empaquetamiento --- es en última instancia un mecanismo para mover el punto de
acuerdo eventual dentro del ZOPA hacia el valor de reserva del que aplica la táctica.

Los resultados de anclaje de \textcite{galinsky2001anchoring} inciden directamente en la
pregunta de quién hace la primera oferta dentro del ZOPA. La alta correlación $r = .93$ entre
la primera oferta y el acuerdo final en condiciones sin defensa significa que la parte que
ancla primero --- con un valor ambicioso pero que respeta el ZOPA --- se ventaja
estructuralmente.
% galinsky2001anchoring: ventaja del ancla de la primera oferta, se aplica dentro del ZOPA
\textcite{northcraft1987anchoring} confirman que incluso los negociadores expertos absorben el
ancla y ajustan de manera insuficiente; la defensa, también identificada por
\textcite{galinsky2001anchoring}, es centrarse exclusivamente en el propio precio objetivo y
el BATNA de la contraparte, en lugar de comprometerse con el ancla como punto de referencia.
% galinsky2001anchoring: enfocarse en el propio objetivo / BATNA de la contraparte anula el efecto del ancla

\begin{example}[Contrato empresarial y una hoja de términos de Series~A]
\textbf{Contrato empresarial.} La SaaS está negociando un contrato empresarial. El BATNA del
vendedor es cerrar tres cuentas adicionales de mercado medio a \money{18000} cada una (total
\money{54000} anuales), con menor sobrecarga de soporte; \(RV_{\text{seller}} = \money{60000}\)
tras neto de costos de transición. El BATNA del comprador empresarial es una construcción
interna personalizada estimada en \money{110000} en ingeniería y mantenimiento;
\(RV_{\text{buyer}} = \money{100000}\) tras un recorte de riesgo del \qty{10}{\percent} por
incertidumbre en la construcción. El ZOPA es \([\money{60000}, \money{100000}]\), con amplitud
\(\Delta = \money{40000}\). El vendedor abre en \money{95000} --- un ancla cercana al techo del
comprador, fundamentada en el valor declarado de la plataforma e implementaciones comparables
nombradas --- y la negociación se cierra en \money{82000}: dentro del ZOPA, por encima de ambos
valores de reserva, y sostenible para ambas partes. El excedente de \money{22000} por encima del
valor de reserva del vendedor y el excedente de \money{18000} por debajo del valor de reserva del
comprador reflejan la ventaja de anclaje del primer movimiento y la disciplina de concesión de
ambas partes.

\textbf{Hoja de términos de Series~A.} La misma lógica se aplica al financiamiento de capital
(\cref{ch:fundraising-and-dilution}). Suponga que el BATNA del fundador es una hoja de términos
competidora con una valoración premoney de \money{12000000}.
\(RV_{\text{founder}}\) es la dilución implícita por esa valoración más el costo de la demora.
El BATNA del inversor líder es desplegar capital en una empresa comparable a una valoración
premoney de \money{15000000}; su \(RV_{\text{investor}}\) es la valoración premoney máxima a la
cual el acuerdo aún cumple su umbral de rendimiento. El ZOPA es la banda de valoración entre
esos dos umbrales. El fundador que ancla primero con una cifra premoney específica respaldada en
un comparable nombrado --- con criterios objetivos, según \textcite{raiffa1982art} --- jala el
acuerdo hacia su valor de reserva. El inversor que retrasa revelar su tasa de umbral de toda la
cartera protege el suyo. Ambos operan dentro del marco de negociación basada en principios.
\end{example}

\begin{admonition}{Advertencia}
Dos modos de falla son comunes. Primero, anclar más allá del valor de reserva de la contraparte
--- abriendo de manera tan agresiva que la otra parte se retira en lugar de contraofertar ---
colapsa el ZOPA al señalar mala fe. El ancla debe ser ambiciosa pero plausible, fundamentada en
criterios objetivos. Segundo, la maldición del ganador: en las subastas competitivas, el ganador
es sistemáticamente la parte que más sobrestimó el valor del activo
\parencite{thaler1988winnerscurse}. Un fundador que «gana» una contratación al licitar muy por
encima del valor de reserva del candidato ha pagado de más en relación con el mercado; una firma
de capital de riesgo que gana un acuerdo al licitar muy por encima de su umbral de rendimiento
ha destruido el rendimiento esperado.
% thaler1988winnerscurse: maldición del ganador --- el postor ganador sobrestima el valor del activo
Ambos errores se remedian anclando el valor de reserva a la utilidad monetaria del BATNA, no
al deseo.
\end{admonition}


\section{Poder e influencia sostenible}\label{sec:power-trust}
% complejidad: 3

Los marcos de negociación de \cref{sec:negotiation} suponen que las partes se reúnen en la
mesa. Antes de la mesa, el poder determina si la reunión ocurre, qué agenda se establece y
cuyo encuadre da forma al ancla de apertura. \textcite{pfeffer2010power} proporciona el análisis
organizacional basado en evidencia: el poder se acumula en quienes controlan las dependencias
de recursos críticos, construyen coaliciones amplias y demuestran competencia de manera visible.
El contraste con la mitología popular del género vale la pena nombrarlo brevemente.
\textcite{greene1998laws} construye una taxonomía de tácticas maquiavélicas a partir de anécdotas
históricas; no se deriva de evidencia sistemática. Las estrategias que endosa --- ocultamiento,
creación estratégica de dependencia --- son precisamente aquellas que destruyen la confianza
necesaria para retener empleados, socios y clientes a lo largo del tiempo. El análisis
estructural de Pfeffer describe dinámicas que se acumulan; las tácticas de manipulación
anecdóticas se erosionan.

El fundamento duradero de la influencia sostenible es la confianza. \textcite{maister2000trusted}
descomponen la confianza en sus componentes estructurales como heurística para profesionales:
\begin{equation}\label{eq:trust}
  T \;=\; \frac{C + R + I}{S},
\end{equation}
donde \(C\) es la credibilidad (la calidad de lo que se dice), \(R\) es la fiabilidad (la
consistencia entre lo que se dice y lo que se hace), \(I\) es la intimidad (el grado en que la
contraparte cree que se manejará la información sensible con discreción) y \(S\) es la
orientación al yo (el grado en que la contraparte percibe el interés propio como centrado en
uno mismo en lugar de en ellos). La estructura de la ecuación es el hallazgo operativo:
\(S\) está en el denominador, lo que significa que incluso una alta credibilidad y fiabilidad
colapsan si la contraparte lee al asesor como interesado en sí mismo.

La heurística para profesionales de \cref{eq:trust} está fundamentada en el modelo organizacional
integrativo de \textcite{mayer1995model}, quienes formalizan la confiabilidad como una estructura
aditiva de tres dimensiones observables de manera independiente: capacidad (competencia específica
del dominio), benevolencia (la percepción de que el fideicomisario desea el éxito del fideicomitente,
no solo el propio) e integridad (consistencia entre palabras y acciones y adhesión percibida a
principios aceptables). \textcite{colquitt2007meta} proporcionan la validación metaanalítica en
132 muestras independientes: la correlación corregida entre cada dimensión y el nivel de
confianza resultante es \(r_c = .67\) para la capacidad, \(r_c = .63\) para la benevolencia y
\(r_c = .62\) para la integridad.
% colquitt2007meta: capacidad r_c=.67, benevolencia r_c=.63, integridad r_c=.62, N~6800–7300 por dimensión
Los tres componentes contribuyen con peso casi igual --- ninguna dimensión domina --- lo que
coincide con la estructura del numerador equilibrado de \cref{eq:trust}. De manera crítica, la
confianza se traduce directamente en el comportamiento de toma de riesgos que hace que un
responsable de adquisiciones firme un contrato plurianual: la correlación corregida entre
confianza y toma de riesgos es \(r_c = .42\) (\(N = 1{,}384\))
\parencite{colquitt2007meta}.
% colquitt2007meta: confianza -> toma de riesgos r_c=.42, N=1384

Reducir la orientación al yo --- hacer preguntas en lugar de presentar argumentos, divulgar los
conflictos de manera proactiva, ocasionalmente dirigir a un cliente a un competidor que sea una
mejor opción --- multiplica el numerador de confianza sin añadirle nada. Esto no es mera
altruismo: un cliente que entra a un contrato con alta confianza (\(T\)) exhibe menor
\emph{churn} (la tasa mensual de \emph{churn} del \qty{0.65}{\percent} asumida en el modelo
de referencia se predica en que los clientes permanezcan convencidos de la benevolencia del
proveedor a través de cada ciclo de renovación), mayor disposición a expandir (moviéndose del
nivel estándar de \money{50}/mes hacia el nivel empresarial) y menor CAC a través de referencias
que se originan dentro de la red de confianza en lugar de la adquisición pagada (\cref{eq:cac}).
La implicación para el ROIC (\cref{eq:roic}) se deriva directamente: la confianza reduce el
numerador del CAC, reduce la attrición impulsada por \emph{churn} en el denominador del CLV y
amplía el margen --- cada uno de los cuales eleva el retorno sobre el capital invertido.

\begin{admonition}{Controvertido}
\Cref{eq:trust} es una heurística para profesionales, no una ecuación de regresión estimada.
\textcite{maister2000trusted} la derivan de la práctica consultora, no de un estudio empírico
controlado. Su virtud es estructural --- identifica correctamente a \(S\) como el divisor
multiplicativo --- pero la razón precisa no está calibrada empíricamente. La ciencia subyacente
es \textcite{mayer1995model} y \textcite{colquitt2007meta}; la ecuación es el mnemónico que
hace operativa la ciencia. Utilice las correlaciones de componentes de \textcite{colquitt2007meta}
para priorizar dónde invertir: como la capacidad, la benevolencia y la integridad contribuyen
casi por igual, un déficit en cualquier dimensión individual no se compensa con un excedente en
las demás.
\end{admonition}

\begin{example}[La confianza del fundador como activo que se acumula]
Considere a un fundador de SaaS que construye relaciones con inversores y talento durante dos
años antes de una ronda de Series~A (\cref{ch:fundraising-and-dilution}). El fundador da cuatro
pasos visibles, cada uno correspondiendo a un componente de \cref{eq:trust}:

\begin{itemize}
  \item \textbf{Credibilidad (\(C\)):} publica actualizaciones mensuales de métricas a una
    pequeña lista de inversores ángel y operadores, incluyendo los meses en que el crecimiento
    estuvo por debajo del objetivo, con atribución honesta de la causa raíz. Esto construye la
    señal de \emph{capacidad} de \textcite{mayer1995model} a través de la competencia demostrada
    en el dominio y la señal de integridad a través de la consistencia entre palabras y acciones.
  \item \textbf{Fiabilidad (\(R\)):} cumple cada compromiso adquirido con los inversores iniciales
    --- fechas de hoja de ruta del producto, hitos de contratación, plazos de respuesta ---
    sin excepción durante ocho trimestres. El canal de integridad $r_c = .62$
    \parencite{colquitt2007meta} se acumula con cada promesa cumplida.
    % colquitt2007meta: integridad r_c=.62, antecedente de confianza
  \item \textbf{Intimidad (\(I\)):} aconseja en confidencialidad a dos candidatos potenciales que
    no se unan a la empresa todavía porque el rol no está bien definido; ambos se unen un año
    después cuando el rol es real. Esto demuestra benevolencia (el canal $r_c = .63$) a un costo
    visible para la velocidad de contratación a corto plazo.
    % colquitt2007meta: benevolencia r_c=.63, antecedente de confianza
  \item \textbf{Orientación al yo (\(S\)):} en las reuniones con inversores, abre con preguntas
    sobre las preocupaciones actuales de la cartera de la firma antes de describir las necesidades
    de la empresa.
\end{itemize}

Las consecuencias medibles aparecen en la ronda. Las referencias provenientes de la red de
confianza cuestan cero en CAC pagado (\cref{eq:cac}). Dos de los tres inversores de Series~A
que el fundador más deseaba toman reuniones sin una presentación cálida, porque el historial de
credibilidad ya es legible. El líder cierra con una valoración premoney en la mitad superior del
rango modelado, en parte porque los inversores con puntuaciones altas de \(T\) exigen un menor
descuento por incertidumbre. Tres de las seis contrataciones clave de ingeniería llegan a través
de referencias de talento del candidato al que el fundador aconsejó anteriormente que se fuera
--- la confianza convirtiéndose en eliminación del CAC de reclutamiento. El mecanismo de
acumulación es el mismo que en \cref{eq:clv}: cada período de comportamiento que refuerza la
confianza eleva el valor presente neto de la relación al reducir la tasa de descuento implícita
que la contraparte aplica a las interacciones futuras.
\end{example}
